\documentclass{leveldoc}
\DocUses{foundations,first_bend,slingshot,around_the_wall,einzel_lens}
\DocFingerprint{57aff31382571059}

\begin{document}

\section{What the level asks}\label{sec:asks}

The game's text:
\leveldescription

One shot leaves $(\val{xLaunch}, \val{yLaunch})$ with charge $q = \val{q}$, mass
$m = \val{m}$ and kinetic energy $T_0 = \val{Tzero}$, aimed along $+x$. The
detector is the rectangle $\val{xDlo} \le x \le \val{xDhi}$,
$\val{yDlo} \le y \le \val{yDhi}$. An entry counts only when the kinetic energy
lies in $[T_{\min}, T_{\max}] = [\val{Tlo}, \val{Thi}]$. There is no direction
cone. The description's quarter of the launch energy is $\val{Tquarter}$; the
file's ceiling is lower than that quarter by $\val{quarterGap}$, and the file
also refuses an arrival below $\val{Tlo}$.

Two electrodes are already on the board, $Q = \val{Qlo}$ at
$(\val{xLo}, \val{yLo})$ and $Q = \val{Qhi}$ at $(\val{xHi}, \val{yHi})$. The
player may add at most $\val{maxCharges}$ further charges, of either sign, inside
the region $\val{xRegLo} \le x \le \val{xRegHi}$,
$\val{yRegLo} \le y \le \val{yRegHi}$, and no magnets. A flight ends when the
particle touches a charge, leaves the world, enters the detector or runs out of
time (\xsecref{foundations}{sec:board}); $t_{\max} = \val{tmax}$.

The reference solution adds $Q = \val{Qref}$ at $(\val{xRef}, \val{yRef})$, the
particle's sign. The added charge and the two electrodes brake the ion, and it
enters the box inside the window without the forward momentum changing sign.
This document integrates that three-charge flight, the two electrodes alone, and
the same node with half the added charge, by the fixed-step method of
\levelref{around_the_wall}, and applies the kinetic window at the entry.

\section{Setup}\label{sec:setup}

The level file, field by field (the appendix reproduces it):
\setuptable

In symbols: $c = \val{c}$, the particle has radius $a = \val{aRad}$, and every
charge has radius $R = \val{R}$. The launch has
\[
  \gamma_0 = 1 + \frac{T_0}{mc^2} = \val{gammaZero}, \qquad
  \abs{\vect{p}_0} = \sqrt{\frac{T_0^2}{c^2} + 2mT_0} = \val{pZero}, \qquad
  v_0 = \frac{\abs{\vect{p}_0}}{\gamma_0 m} = \val{vZero},
\]
at $\val{angLaunch}^\circ$ to $+x$. The world is $-1 \le x \le \val{xWorld}$,
$-1 \le y \le \val{yWorld}$. The electrodes are equal and sit symmetrically about
$y = \val{yMid}$, which is the launch height. Their couplings are
$\kappa_\mathrm{l} = qQ_\mathrm{l} = \val{kLo}$ and
$\kappa_\mathrm{h} = qQ_\mathrm{h} = \val{kHi}$. The reference adds
$\kappa = \val{kRef}$ at $(\val{xRef}, \val{yRef})$, so the reference flight has
$N = \val{nRef}$ centres. The electrodes alone have $N = \val{nPair}$. The third
placement puts $\kappa = \val{kHalf}$ at the same node $(\val{xHalf}, \val{yHalf})$,
and has $N = \val{nHalf}$.

\begin{figure}[htbp]
\centering
\begin{tikzpicture}
\begin{axis}[axis equal image, width=0.88\linewidth, xmin=-1.8, xmax=32.4, ymin=-1.8, ymax=21.8,
  xtick={0,5,...,30}, ytick={0,5,...,20}, tick label style={font=\small},
  xlabel={$x$}, ylabel={$y$}, label style={font=\small}]
  \draw[black!45, dashed] (-1,-1) rectangle (31,21);
  \draw[black!50] (0,0) rectangle (30,20);
  \draw[black!40, dashed] (\rawval{xRegLo},\rawval{yRegLo}) rectangle (\rawval{xRegHi},\rawval{yRegHi});
  \fill[green!50!black, fill opacity=0.16] (\rawval{xDlo},\rawval{yDlo}) rectangle (\rawval{xDhi},\rawval{yDhi});
  \draw[green!40!black] (\rawval{xDlo},\rawval{yDlo}) rectangle (\rawval{xDhi},\rawval{yDhi});
  \fill[blue!70!black, fill opacity=0.35] (\rawval{xLo},\rawval{yLo}) circle[radius=\rawval{R}];
  \draw[blue!70!black] (\rawval{xLo},\rawval{yLo}) circle[radius=\rawval{R}];
  \fill[blue!70!black, fill opacity=0.35] (\rawval{xHi},\rawval{yHi}) circle[radius=\rawval{R}];
  \draw[blue!70!black] (\rawval{xHi},\rawval{yHi}) circle[radius=\rawval{R}];
  \fill[red!70!black, fill opacity=0.35] (\rawval{xRef},\rawval{yRef}) circle[radius=\rawval{R}];
  \draw[red!70!black] (\rawval{xRef},\rawval{yRef}) circle[radius=\rawval{R}];
  \addplot[thick, purple!70!black] table {pathpair.dat};
  \addplot[thick, black!55] table {pathhalf.dat};
  \addplot[very thick, blue!70!black] table {pathref.dat};
  \draw[-{Stealth[length=6pt]}, very thick, blue!70!black] (\rawval{xLaunch},\rawval{yLaunch}) --
    (\rawval{arrowX},\rawval{arrowY});
  \fill (\rawval{xLaunch},\rawval{yLaunch}) circle[radius=0.11];
  \node[font=\scriptsize, blue!50!black, anchor=north] at (\rawval{xLo},\rawval{yLo}) {$Q_\mathrm{l}$};
  \node[font=\scriptsize, blue!50!black, anchor=south] at (\rawval{xHi},\rawval{yHi}) {$Q_\mathrm{h}$};
  \node[font=\scriptsize, red!60!black, anchor=north] at (\rawval{xRef},\rawval{yRef}) {$Q$};
  \node[font=\scriptsize, green!30!black, anchor=west] at (30.15,10) {detector};
\end{axis}
\end{tikzpicture}
\caption{Blue is the reference flight: it rises, turns down, and enters the left
face of the detector. Purple is the two electrodes alone, straight along the
launch line into that same face. Grey is half the added charge, at the same red
node: it passes above the detector and leaves through the right edge of the
world. The blue discs are the electrodes and the red disc is the added charge.
The dashed rectangle is the region in which the player may place charges.}
\label{fig:board}
\end{figure}

\section{Principles and model}\label{sec:model}

\begin{principles}
  \item The field of a charged sphere and the force on a non-polarizable particle,
    \xeqref{foundations}{eq:coulomb} and \xeqref{foundations}{eq:sphereforce}.
  \item The relativistic equation of motion \xeqref{foundations}{eq:eom}, the launch
    \xeqref{foundations}{eq:launch} and the energy \xeqref{foundations}{eq:energy}. The
    motion stays in the plane (\xsecref{foundations}{sec:plane}).
  \item The board's rules (\xsecref{foundations}{sec:board}). The first entry into the
    detector decides. This level adds a kinetic window and no direction cone.
  \item The classical Runge--Kutta method \xeqref{foundations}{eq:rk4} and the fixed-step
    flight of \xsecref{around_the_wall}{sec:method}. With two or more centres the orbit
    equation does not separate (\xeqref{slingshot}{eq:noseparation}).
  \item The kinetic energy \xeqref{einzel_lens}{eq:kinetic}. The window below is the
    condition that this energy, at the entry, lie between the file's two bounds.
  \item The radiated power \xeqref{foundations}{eq:larmor}, with the acceleration
    \xeqref{first_bend}{eq:acceleration}, summed by Simpson's rule
    \xeqref{foundations}{eq:simpson}.
\end{principles}

The level has no magnets, so $\vect{B} = \vect{0}$. The force is
\xeqref{around_the_wall}{eq:force} with $N = \val{nRef}$ on the reference, the
centres in the file's order: the lower electrode, the upper electrode, then the
added charge. The state $(x, y, p_x, p_y)$ obeys \xeqref{foundations}{eq:eom}
with that force until the first event. The potential $V$ in
\xeqref{around_the_wall}{eq:energy} is the same superposition, and the kinetic
energy at a point is \xeqref{einzel_lens}{eq:kinetic}.

\begin{assumption}{The charges are held in place}\label{as:fixed}
None of the charges recoils. That is a rule of the game
(\xsecref{foundations}{sec:radiation}).
\end{assumption}

\begin{assumption}{No radiation}\label{as:radiation}
The radiated energy is not taken from the particle. Over the three flights of this
document the largest fraction of $T_0$ is $\val{radFracMax}$ (§\ref{sec:hand}). The
game makes the same approximation and reports the same integral (\physref{8}).
\end{assumption}

\begin{assumption}{A point particle}\label{as:point}
$a = \val{aRad}$, and the particle does not polarize, so
\xeqref{foundations}{eq:sphereforce} is exact. The closest approach on the
reference flight is $\val{RefRclose}$, to the \val{RefRname}, a clearance of
$\val{RefClear}$ outside the sphere. The electrodes alone stay $\val{PairClear}$
clear, and the half-charge flight stays $\val{HalfClear}$ clear. The samples that
locate these distances are at most $\val{sampleArc}$ apart (§\ref{sec:method}),
far below the clearances.
\end{assumption}

\begin{assumption}{The step resolves the landing}\label{as:step}
The reported step is $h = \val{H}$. Halving it estimates the largest error of an
end state as $\val{errMax}$ (§\ref{sec:hand}). That is far below the height of the
detector and far below the gap between the reference arrival energy and either
edge of the window. The tolerance used in §\ref{sec:comparison} is the larger of
$10^{-9}$ and ten times this estimate.
\end{assumption}

\section{Derivation}\label{sec:derivation}

\subsection{The kinetic window}\label{sec:window}

The launch fixes $\gamma_0$ by \xeqref{foundations}{eq:launch}, so the kinetic
energy carried by a momentum is
\begin{equation}\label{eq:arrival}
  T(\vect{p}) = (\gamma - 1)\,mc^2, \qquad
  \gamma = \sqrt{1 + \frac{\abs{\vect{p}}^2}{m^2 c^2}}.
\end{equation}
While \xeqref{around_the_wall}{eq:energy} holds, this is the same number as
\xeqref{einzel_lens}{eq:kinetic}. Substituting $W = \gamma_0 mc^2 + V(\vect{x}_0)$
and $\gamma_0 mc^2 = mc^2 + T_0$ gives
\begin{equation}\label{eq:from-potential}
  T(\vect{x}) = T_0 + V(\vect{x}_0) - V(\vect{x}).
\end{equation}
The file accepts an entry only when
\begin{equation}\label{eq:window}
  T_{\min} \le T \le T_{\max},
\end{equation}
with $T_{\min} = \val{Tlo}$ and $T_{\max} = \val{Thi}$. For the potential
expression that is a bound on the rise of $V$ between the launch and the entry,
\begin{equation}\label{eq:rise}
  T_0 - T_{\max} \le V(\vect{x}_\mathrm{e}) - V(\vect{x}_0) \le T_0 - T_{\min}.
\end{equation}
The lower bound is $\val{riseMin}$ and the upper bound is $\val{riseMax}$. The
ion must be slowed by at least $\val{riseMin}$, and by no more than $\val{riseMax}$.
The description's quarter, $\val{Tquarter}$, would have asked only for a rise of
$T_0 - \val{Tquarter}$. The file asks for $\val{quarterGap}$ more, and it also
rejects a landing below $\val{Tlo}$.

The test itself uses \eqref{eq:arrival} at the entry momentum, not
\eqref{eq:from-potential}. The signed margin is
\begin{equation}\label{eq:margin}
  \mu = \frac{\min\bigl(T - T_{\min},\, T_{\max} - T\bigr)}{\abs{T_{\max}}}.
\end{equation}
The entry counts when $\mu \ge 0$, so both endpoints belong to the window. When
$\mu < 0$ the outcome is Rejected and the reported state stays the entry
(\physref{6.1}). A flight that never enters the box is not tested. On the flights
below, \eqref{eq:arrival} and \eqref{eq:from-potential} differ by less than the
drift of $W$ (§\ref{sec:hand}), so the rise \eqref{eq:rise} decides the same
question as the margin.

\subsection{Why the flight is integrated}\label{sec:noseparation}

With one charge, angular momentum closes the motion into the orbit of
\levelref{first_bend}. Two fixed centres already fail to separate in the
relativistic problem: the term \xeqref{slingshot}{eq:noseparation} does not split.
Three centres, two of them symmetric and the third off the axis, give that
obstruction nothing that would remove it. Each flight is the initial-value
problem for \xeqref{around_the_wall}{eq:force}, and §\ref{sec:method} computes it
to a stated accuracy.

\subsection{The method}\label{sec:method}

The state is advanced by the method of \xsecref{around_the_wall}{sec:method}: the
classical Runge--Kutta method \xeqref{foundations}{eq:rk4} at the fixed step
$h = \val{H}$, the crossing on the cubic Hermite interpolant
\xeqref{slingshot}{eq:hermite}, $\val{nSample}$ samples and $\val{nBisect}$
bisections of that interpolant, and the enlarged-box test. The position reported
on a face has its face coordinate set to the face value. The momentum is one
further step of the method, of the length up to the crossing, and
\eqref{eq:window} is applied to \eqref{eq:arrival} for that momentum. The error of
the reported step is a fifteenth of its difference from the step $2h$
(\xsecref{foundations}{sec:numerics}).

On each step the same interpolant is sampled at nine equally spaced fractions,
and only up to the crossing when the step contains the event. Those samples give
the kinetic-energy minimum, the smallest forward momentum, and the closest
approach. A point between them lies within $\val{sampleArc}$ of a sample, so
those three quantities are located to that tolerance. The end state is held to
the step-halving estimate.

\section{The reference solution by hand}\label{sec:hand}

\begin{procedure}
  \item\label{st:launch} \emph{The launch} (\xeqref{foundations}{eq:launch}):
    $\gamma_0 = \val{gammaZero}$ and $\abs{\vect{p}_0} = \val{pZero}$, aimed at
    $\val{angLaunch}^\circ$. The force is $(\val{RefFx}, \val{RefFy})$. The
    horizontal part is negative: every charge stands ahead of the particle and
    repels. The vertical part is positive. The added charge sits below the beam,
    so it pushes the particle up; the electrodes cancel on the axis.
  \item\label{st:energy} \emph{The energy}: $V(\vect{x}_0) = \val{RefVzero}$,
    $W = \val{RefW}$ and $W - mc^2 = \val{RefWbind} > 0$. The world's edge is
    energetically open. \eqref{eq:rise} asks the potential to rise by at least
    $\val{riseMin}$ and by at most $\val{riseMax}$.
  \item\label{st:flight} \emph{The flight}: the method of §\ref{sec:method} with
    $h = \val{H}$. The states every five time units, and the drift of the energy:
    \begin{center}\small\val{states}\end{center}
    The path climbs to $y = \val{RefYmax}$ and then comes down. The kinetic energy
    falls to $\val{RefTmin}$ at $t = \val{RefTTmin}$, at
    $(\val{RefXTmin}, \val{RefYTmin})$, near the upper electrode. The smallest
    forward momentum on the samples is $\val{RefPxMin}$, at
    $t = \val{RefTPxMin}$. It stays positive, so the particle does not turn back.
  \item\label{st:entry} \emph{The entry.} The step from $t_n = \val{RefTbefore}$,
    at $(\val{RefXbefore}, \val{RefYbefore})$, crosses the \val{RefFace} face of
    the detector at the fraction $u = \val{RefU}$. The residual before the face
    coordinate is assigned is $\val{RefFaceRes}$, so $x_\mathrm{e} = \val{RefX}$,
    $y_\mathrm{e} = \val{RefY}$ and $t_\mathrm{e} = \val{RefT}$, after
    $\val{RefSteps}$ steps. One step of length $t_\mathrm{e} - t_n$ gives
    $\vect{p}_\mathrm{e} = (\val{RefPx}, \val{RefPy})$, at $\val{RefAng}^\circ$
    to $+x$. The ordinate lies between $\val{yDlo}$ and $\val{yDhi}$.
  \item\label{st:window} \emph{The window} \eqref{eq:window}. From the entry
    momentum, \eqref{eq:arrival} gives $T = \val{RefTarr}$. The potential
    expression \eqref{eq:from-potential} gives $\val{RefTarrV}$. They differ by
    $\val{RefTgap}$. The rise of $V$ is $\val{RefRise}$, which lies between
    $\val{riseMin}$ and $\val{riseMax}$. The margin \eqref{eq:margin} is
    $\val{RefMargin} > 0$: the arrival sits $\val{RefRoomLo}$ above the floor and
    $\val{RefRoomHi}$ below the ceiling. The entry counts.
  \item\label{st:check} \emph{The checks}: the energy stays within $\val{RefWdev}$
    of $W$. Halving the step estimates the end state as $\val{RefErr}$. The same
    outcome appears at $h$ and at $2h$. The closest approach, $\val{RefRclose}$ to
    the \val{RefRname}, leaves a clearance $\val{RefClear}$.
  \item\label{st:rad} \emph{The radiation}: the power \xeqref{foundations}{eq:larmor}
    with the acceleration \xeqref{first_bend}{eq:acceleration}, summed by Simpson's
    rule on the steps of $h/2$. The integral is $\val{RefRad}$, the fraction
    $\val{RefRadFrac}$ of $T_0$. The relative difference from the same sum at step
    $h$, divided by fifteen, is $\val{RefRadErr}$.
\end{procedure}

\begin{result}[A landing inside the window]
The reference enters the detector through its left face at
$(\val{RefX}, \val{RefY})$, at $t = \val{RefT}$, with kinetic energy
$\val{RefTarr}$ inside $[\val{Tlo}, \val{Thi}]$, heading $\val{RefAng}^\circ$
to $+x$. The forward momentum has stayed at least $\val{RefPxMin}$. The level
is solved.
\end{result}

By hand, Step~\ref{st:flight} is $\val{RefSteps}$ steps of four evaluations of
the force. Across the three flights of this document the largest step-halving
estimate is $\val{errMax}$.

\section{Comparison with the game}\label{sec:comparison}

\comparisontable

The three outcomes agree, and so do the three end positions, times and momenta.
The arrival kinetic energy is compared on the two flights that enter the box,
with the state tolerance, because \eqref{eq:arrival} is computed from the entry
momentum. The tolerance on a state is $\val{tolFlight}$. Ten times the largest
step-halving estimate $\val{errMax}$ lies below $10^{-9}$, so the tolerance is
that floor. The three radiated energies are compared relatively, with tolerance
$10^{-4}$. The state tolerance and the radiation tolerance were fixed from these
estimates before the comparison.

The largest absolute difference is the reference time, $\val{cmpRefTDiff}$. The
reference ordinate differs by $\val{cmpRefYDiff}$. The difference is many times
$\val{errMax}$ and a fraction of the tolerance. The game's verified end states are
good to about $10^{-10}$ (\xsecref{foundations}{sec:verification}), so a
difference of this size is what that accuracy allows when three centres are
summed.

The reference abscissa agrees exactly, as do the electrodes-only ordinate and its
vertical momentum. The quadrature uncertainties are $\val{RefRadErr}$,
$\val{PairRadErr}$ and $\val{HalfRadErr}$, all inside the relative tolerance. The
largest relative difference is the electrodes-only flight, $\val{cmpPairRadDiff}$.

\section{Other solutions}\label{sec:others}

The electrodes are not enough, and half the reference charge misses the box. The
method of §\ref{sec:method} integrates both (\code{placements.json}), each at
$h = \val{H}$ and at $2h$.
\begin{enumerate}[label=\arabic*.]
  \item The two electrodes alone. They are equal and symmetric about
    $y = \val{yMid}$, so the launch force is $(\val{PairFx}, \val{PairFy})$ and
    the ordinate never leaves $\val{yLaunch}$: the deviation is $\val{PairYdev}$.
    The particle enters the \val{PairFace} face at $(\val{PairX}, \val{PairY})$,
    at $t = \val{PairT}$, with momentum $(\val{PairPx}, \val{PairPy})$, at
    $\val{PairAng}^\circ$.

    The rise of $V$ is $\val{PairRise}$, short of the $\val{riseMin}$ that
    \eqref{eq:rise} requires. The arrival energy is $\val{PairTarr}$, above
    $\val{Thi}$. The margin is $\val{PairMargin} < 0$, so the outcome is Rejected
    and the reported state stays the entry. The forward momentum stays at least
    $\val{PairPxMin}$. The kinetic energy never falls below $\val{PairTmin}$,
    which is itself above the ceiling. The estimate is $\val{PairErr}$. The
    radiated energy is $\val{PairRad}$.
  \item Half the added charge, $Q = \val{Qhalf}$ at $(\val{xHalf}, \val{yHalf})$.
    The path rises to $\val{HalfYmax}$. At the detector's left edge,
    $x = \val{xDlo}$, the ordinate is $\val{HalfYDlo}$, which is
    $\val{HalfAboveDlo}$ above the top of the box. At the right edge,
    $x = \val{xDhi}$, it is $\val{HalfYDhi}$, still $\val{HalfAboveDhi}$ above
    the top. The particle never enters the box. It leaves through the
    \val{HalfFace} face of the world at $(\val{HalfX}, \val{HalfY})$, at
    $t = \val{HalfT}$, with momentum $(\val{HalfPx}, \val{HalfPy})$, at
    $\val{HalfAng}^\circ$.

    The window is not tested. Along the whole flight the kinetic energy stays at
    least $\val{HalfTmin}$, above $\val{Thi}$, so an entry anywhere would have
    failed \eqref{eq:window} as well. The forward momentum stays at least
    $\val{HalfPxMin}$. The estimate is $\val{HalfErr}$. The radiated energy is
    $\val{HalfRad}$.
\end{enumerate}

The reference needs the full added charge. The electrodes alone brake the ion,
and not enough: it arrives on the target with $\val{PairTarr}$, above
$\val{Thi}$. Half the added charge lifts the path just over the box and never lands.
The full charge lifts it further, the upper electrode turns it back down into the
box, and the three centres together produce the rise $\val{RefRise}$.
Figure~\ref{fig:board} shows the three paths. Section~\ref{sec:comparison}
records the outcomes.

\section{What the game computes}\label{sec:engine}

The game integrates \xeqref{foundations}{eq:eom} with its eighth-order Runge--Kutta
method at the tolerances $10^{-10}$ and $10^{-12}$ (\physref{5}). It finds the
crossing on its interpolant (\physref{6}) and applies \eqref{eq:window} to
\eqref{eq:arrival} at that entry (\physref{6.1}). It verifies the outcome by
comparing its two runs (\physref{7}). The model is this document's, including
\ref{as:radiation} (\physref{8}). The agreement in §\ref{sec:comparison} is
between the fixed-step flight of §\ref{sec:method}, the window of
§\ref{sec:window}, and that integration.

\section{Further reading}\label{sec:reading}

\begin{itemize}
  \item \levelref{einzel_lens}: the kinetic energy \xeqref{einzel_lens}{eq:kinetic},
    from which the window \eqref{eq:window} is derived here.
  \item \levelref{around_the_wall}: the force of several centres, the conserved
    energy, and the fixed-step method, used here at $N = \val{nRef}$.
  \item \levelref{slingshot}: why two fixed centres have no closed relativistic
    orbit, and the Hermite crossing.
  \item \levelref{first_bend}: the acceleration that enters the radiated power.
  \item \xsecref{foundations}{sec:board} and \xsecref{foundations}{sec:numerics}:
    the order of events, and the step-halving estimate.
\end{itemize}

\levelfileappendix

\end{document}
